% SEGMENT OLP-0029-S01
% Part:sets-functions-relations
% Chapter: size-of-sets
% Section: enumerations

\documentclass[../../../include/open-logic-section]{subfiles}

\begin{document}

\olfileid{sfr}{siz}{enm}

\olsection{பட்டியலாக்கங்களும் \usetoken{S}{enumerable} கணங்களும்}

\begin{editorial}
இந்தப் பிரிவு கணங்களின் பட்டியலாக்கங்களை,
$\PosInt$ இலிருந்து செல்லும் மேற்கோர்த்தல் சார்புகளாக வரையறுத்து
விவாதிக்கிறது. கணிதப் பின்னணி குறைவாக உள்ள வாசகர்களுக்காக
விளக்கம் படிப்படியாகச் செல்கிறது.
சுருக்கமான மாற்று வடிவம் \olref[enm-alt]{sec} இல் கொடுக்கப்பட்டுள்ளது.
அதில் பட்டியலாக்கங்கள் வேறு விதமாக,
$\Nat$ உடனான (அல்லது அதன் ஒரு தொடக்கத் துண்டுடனான)
இருபுறச் சார்புகளாக வரையறுக்கப்படுகின்றன.
\end{editorial}

% SEGMENT OLP-0029-S02
\begin{explain}
கணங்களின் !!{element}s[acc] பட்டியலிட்டு அவற்றிற்கு ஏற்கெனவே
எடுத்துக்காட்டுகள் கொடுத்துள்ளோம்.
ஒரு கணம் முடிவுறாததாக இருந்தால்கூட, அதன் !!{element}s[acc]
எவ்வாறு, எப்போது பட்டியலிடலாம் என்பதை மேலும் பொதுவாக விவாதிப்போம்.
\end{explain}

% SEGMENT OLP-0029-S03
\begin{defn}[பட்டியலாக்கம் — முறைசாரா விளக்கம்]
முறைசாரா விளக்கத்தின்படி, ஒரு கணம் $A$ இன் \emph{பட்டியலாக்கம்}
என்பது, $A$ இன் !!{element}s கொண்ட ஒரு பட்டியலாகும்
(அது முடிவுறாததாகவும் இருக்கலாம்);
அதில் $A$ இன் ஒவ்வோர் !!{element}[um] ஏதாவது ஒரு முடிவுறு
இடத்தில் இடம்பெற வேண்டும்.
$A$ க்கு ஒரு பட்டியலாக்கம் இருந்தால்,
$A$ என்பது \emph{!!{enumerable}} எனப்படும்.
\end{defn}

% SEGMENT OLP-0029-S04
\begin{explain}
பட்டியலாக்கங்களைப் பற்றிய சில குறிப்புகள்:
\begin{enumerate}
\item ஒரு தொடக்கத்தைக் கொண்டதும், முதலாவதைத் தவிர
  ஒவ்வோர் !!{element}[um] தனக்கு உடனே முன்பாக ஒரே ஓர்
  !!{element}[acc]க் கொண்டதுமான பட்டியல்களையே
  பட்டியலாக்கங்களாகக் கருதுகிறோம்.
  வேறு சொற்களில், பட்டியலின் முதல் !!{element} க்கும்
  வேறு எந்த !!{element} க்கும் இடையில் முடிவுறு எண்ணிக்கையிலான
  உறுப்புகள் மட்டுமே உள்ளன.
  குறிப்பாக, ஒரு பட்டியலாக்கத்தின் ஒவ்வோர் !!{element}[um]
  முடிவுறு இடத்தைக் கொண்டுள்ளது: முதல் !!{element} இன் இடம் $1$,
  இரண்டாவதன் இடம் $2$, இவ்வாறே தொடர்கிறது.
\item ஒரே கணம் $A$ க்கு, !!{element}s இடம்பெறும் வரிசையில்
  வேறுபடும் வெவ்வேறு பட்டியலாக்கங்கள் இருக்கலாம்:
  $4$, $1$, $25$, $16$, $9$ என்பது முதல் ஐந்து வர்க்க எண்களின்
  கணத்தைப் பட்டியலாக்குவது போலவே,
  $1$, $4$, $9$, $16$, $25$ என்பதும் பட்டியலாக்குகிறது.
\item மீள இடம்பெறுதல்கள் உள்ள பட்டியலாக்கங்களும் பட்டியலாக்கங்களே:
  $1$, $1$, $2$, $2$, $3$, $3$, \dots{} என்பது,
  $1$, $2$, $3$, \dots{} பட்டியலாக்கும் அதே கணத்தைப் பட்டியலாக்குகிறது.
\item ஒரு பட்டியலாக்கத்தைக் குறிப்பிடும்போது வரிசையும்
  மீள இடம்பெறுதலும் \emph{முக்கியமானவை}.
  நேர்ம முழுக்களை $1$, $2$, $3$, $1$, \dots{} எனத் தொடங்கிப்
  பட்டியலாக்கலாம். ஆனால் வழக்கமான முறையில்
  $1$, $2$, $3$, $4$, \dots எனப் பட்டியலாக்கும்போது
  அதன் அமைப்பை எளிதில் காணலாம்.
\item பட்டியலாக்கங்களுக்கு ஒரு தொடக்கம் இருக்க வேண்டும்:
  \dots, $3$, $2$, $1$ என்பது நேர்ம முழுக்களின் பட்டியலாக்கமல்ல;
  ஏனெனில் அதற்கு முதல் !!{element} இல்லை.
  முறைசாரா வரையறையிலிருந்து இது எவ்வாறு பின்தொடர்கிறது
  என்பதைப் புரிந்துகொள்ள, “இந்தப் பட்டியலில் 76 என்ற எண்
  எந்த இடத்தில் வருகிறது?” என்று உங்களையே கேட்டுக்கொள்ளுங்கள்.
\item பின்வருவது நேர்ம முழுக்களின் பட்டியலாக்கமல்ல:
  $1$, $3$, $5$, \dots, $2$, $4$, $6$, \dots\@
  இங்குள்ள சிக்கல் என்னவெனில், இரட்டை எண்கள் முடிவுறு இடங்களில்
  இடம்பெறாமல், $\infty + 1$, $\infty + 2$, $\infty +
  3$ ஆகிய இடங்களில் இடம்பெறுகின்றன.
\item வெற்றுக் கணம் எண்ணிடத்தக்கது:
  அது வெற்றுப் பட்டியலால் பட்டியலாக்கப்படுகிறது!{}
\end{enumerate}
\end{explain}

% SEGMENT OLP-0029-S05
\begin{prop}
$A$ க்கு ஒரு பட்டியலாக்கம் இருந்தால்,
மீள இடம்பெறுதல்கள் இல்லாத ஒரு பட்டியலாக்கமும் அதற்கு உள்ளது.
\end{prop}

% SEGMENT OLP-0029-S06
\begin{proof}
$A$ க்கு $x_1$, $x_2$, \dots{} என்ற ஒரு பட்டியலாக்கம் இருப்பதாகவும்,
அதில் ஒவ்வொரு $x_i$ உம் $A$ இன் ஓர் !!{element} எனவும் கொள்வோம்.
மீண்டும் இடம்பெறும் !!{element}s[acc] நீக்குவதன் மூலம்,
ஒரு பட்டியலாக்கத்திலிருந்து மீள இடம்பெறுதல்களை நீக்கலாம்.
எடுத்துக்காட்டாக, $x_i$ என்பது $x_1$, \dots, $x_{i-1}$ ஆகியவற்றில்
இல்லாத $A$ இன் !!a{element} எனில், அதைப் பட்டியலில்
இடம்பெறச் செய்யலாம். அது ஏற்கெனவே $x_1$, \dots, $x_{i-1}$
ஆகியவற்றில் இடம்பெற்றிருந்தால், $x_i$ ஐப் பட்டியலிலிருந்து நீக்கலாம்.
இவ்வாறு புதிய பட்டியலாக்கம் ஒன்றைப் பெறலாம்.
\end{proof}

பட்டியலாக்கங்களையும் !!{enumerable} கணங்களையும் நன்கு புரிந்துகொண்டு,
அவற்றைப் பற்றிய முடிவுகளை நிறுவுவதற்கு மேலும் துல்லியமான
வரையறை தேவை என்பதை முந்தைய வாதம் காட்டுகிறது.
பின்வரும் வரையறை அதனை வழங்குகிறது.

% SEGMENT OLP-0029-S07
\begin{defn}[பட்டியலாக்கம் — முறையான வரையறை]
$A \neq \emptyset$ என்ற கணத்தின் \emph{பட்டியலாக்கம்} என்பது,
!!{surjective} பண்புடைய ஏதேனும் ஒரு
$f \colon \PosInt \to A$ என்ற சார்பாகும்.
\end{defn}

% SEGMENT OLP-0029-S08
\begin{explain}
முறையான வரையறையும், முடிவுறாததாகவும் இருக்கக்கூடிய பட்டியலைப்
பயன்படுத்தும் முறைசாரா வரையறையும் சமானமானவை என்பதைப் பார்ப்போம்.
முதலில், $\PosInt$ இலிருந்து ஒரு கணம் $A$ க்குச் செல்லும்
எந்த !!{surjective} சார்பும் $A$ ஐப் பட்டியலாக்குகிறது.
அத்தகைய சார்பு, மேலே முறைசாரா வகையில் வரையறுத்தபடி,
$f(1)$, $f(2)$, $f(3)$, \dots என்ற பட்டியலாக்கத்தை நிர்ணயிக்கிறது.
$f$ என்பது !!{surjective} பண்புடையதால்,
$A$ இன் ஒவ்வோர் !!{element}[um] ஏதேனும் ஒரு
$n \in \PosInt$ க்கு $f(n)$ இன் மதிப்பாக இருப்பது உறுதி.
எனவே $A$ இன் ஒவ்வோர் !!{element}[um] பட்டியலில் ஏதாவது ஒரு
முடிவுறு இடத்தில் இடம்பெறுகிறது.
சார்பு !!{injective} பண்புடையதாக இல்லாமல் இருக்கக்கூடும் என்பதால்,
பட்டியலில் மீள இடம்பெறுதல்கள் இருக்கலாம்.
ஆனால் மேலே குறிப்பிட்டபடி அது ஏற்கத்தக்கதே.

மறுபுறம், $A$ இன் எல்லா !!{element}s[acc]யும் பட்டியலாக்கும்
ஒரு பட்டியல் தரப்பட்டால், பின்வருமாறு
!!a{surjective} சார்பு $f\colon \PosInt \to A$ ஐ வரையறுக்கலாம்:
$f(n)$ என்பதைப் பட்டியலின் $n$ ஆவது !!{element} ஆகவும்,
$n$ ஆவது !!{element} இல்லையெனில் பட்டியலின் இறுதி
!!{element} ஆகவும் அமைக்கலாம்.
இது !!a{surjective} சார்பை உருவாக்காத ஒரே நிலை,
$A$ வெற்றுக் கணமாகவும் அதன் காரணமாகப் பட்டியல் வெற்றாகவும்
இருக்கும் நிலையே.
எனவே ஒவ்வொரு வெற்றற்ற பட்டியலும்
!!a{surjective} சார்பு $f\colon \PosInt \to A$ ஐ நிர்ணயிக்கிறது.
\end{explain}

% SEGMENT OLP-0029-S09
\begin{defn}
\ollabel{defn:enumerable}
ஒரு கணம் $A$ வெற்றுக் கணமாக இருக்கும்போது அல்லது அதற்கு
ஒரு பட்டியலாக்கம் இருக்கும்போது, அப்போதும் அப்போது மட்டுமே,
அது !!{enumerable} ஆகும்.
\end{defn}

% SEGMENT OLP-0029-S10
\begin{ex}
நேர்ம முழுக்களை ($\PosInt$) பட்டியலாக்கும் ஒரு சார்பு,
$f(n) = n$ எனக் கொடுக்கப்படும் ஒன்றாதல் சார்பே ஆகும்.
இயல் எண்கள் $\Nat$ ஐப் பட்டியலாக்கும் ஒரு சார்பு,
$g(n) = n - 1$ என்ற சார்பாகும்.
\end{ex}

% SEGMENT OLP-0029-S11
\begin{ex}
பின்வருமாறு கொடுக்கப்படும்
$f\colon \PosInt \to \PosInt$, $g \colon \PosInt \to
\PosInt$ ஆகிய சார்புகள்:
\begin{align*}
f(n) & = 2n \text{ மற்றும்}\\
g(n) & = 2n - 1
\end{align*}
முறையே இரட்டை நேர்ம முழுக்களையும் ஒற்றை நேர்ம முழுக்களையும்
பட்டியலாக்குகின்றன.
ஆயினும், இந்த இரு சார்புகளில் எதுவும் $\PosInt$ இன்
பட்டியலாக்கமல்ல; ஏனெனில் எதுவும் !!{surjective} பண்புடையதல்ல.
\end{ex}

% SEGMENT OLP-0029-S12
\begin{prob}
நேர்ம வர்க்க எண்கள் $1$, $4$, $9$, $16$, \dots ஆகியவற்றின்
ஒரு பட்டியலாக்கத்தை வரையறுக்கவும்.
\end{prob}

% SEGMENT OLP-0029-S13
\begin{ex}
$f(n) = (-1)^{n} \lceil \frac{(n-1)}{2}\rceil$ என்ற சார்பு,
முழுக்களின் கணம் $\Int$ ஐப் பட்டியலாக்குகிறது.
இங்கு $\lceil x \rceil$ என்பது \emph{மேல்முழுச்} சார்பைக் குறிக்கிறது:
அது $x$ ஐ, அதற்குச் சமமாகவோ பெரியதாகவோ உள்ள மிக அருகிலுள்ள
முழுவுக்கு மேல்நோக்கி முழுமைப்படுத்துகிறது.
$f$ ஆனது நேர்ம, எதிர்ம முழுக்களுக்கு இடையில் முன்னும் பின்னும்
“தாவி”, $\Int$ இன் மதிப்புகளை எவ்வாறு உருவாக்குகிறது என்பதைக் கவனிக்கவும்:
% SOURCE CORRECTION OLSIZ-001; see translation/size-notes.tex.
\[
\begin{array}{c c c c c c c c}
f(1) & f(2) & f(3) & f(4) & f(5) & f(6) & f(7) & \dots \\ \\
- \lceil \tfrac{0}{2} \rceil & \lceil \tfrac{1}{2}\rceil & - \lceil \tfrac{2}{2} \rceil & \lceil \tfrac{3}{2} \rceil & - \lceil \tfrac{4}{2} \rceil  & \lceil \tfrac{5}{2}
\rceil & - \lceil \tfrac{6}{2} \rceil & \dots \\ \\
0 & 1 & -1 & 2 & -2 & 3 & -3 & \dots
\end{array}
\]
$f$ பின்வருமாறு வகைப்பிரிவுகளால் வரையறுக்கப்பட்டதாகவும் கருதலாம்:
\[
f(n) = \begin{cases}
  0 & \text{$n = 1$ எனில்}\\
  n/2 & \text{$n$ இரட்டை எனில்}\\
  -(n-1)/2 & \text{$n$ ஒற்றையாகவும் $>1$ ஆகவும் எனில்}
  \end{cases}
\]
\end{ex}

% SEGMENT OLP-0029-S14
\begin{prob}
$A$, $B$ ஆகியவை !!{enumerable} எனில்,
$A \cup B$ என்பதும் எண்ணிடத்தக்கது எனக் காட்டுக.
இதற்கு, !!{surjective} சார்புகள் $f\colon \PosInt \to
A$, $g\colon \PosInt \to B$ உள்ளன எனக் கொண்டு,
!!a{surjective} சார்பு $h\colon \PosInt \to A \cup B$ ஐ
வரையறுத்து, அது !!{surjective} பண்புடையது என நிறுவுக.
$A$ வெற்றுக் கணமாக இருக்கும் நிலையையும்,
$B = \emptyset$ என்ற நிலையையும் கருதுக.
\end{prob}

% SEGMENT OLP-0029-S15
\begin{prob}
$B \subseteq A$ எனவும் $A$ என்பது !!{enumerable} எனவும் இருந்தால்,
$B$ உம் எண்ணிடத்தக்கது எனக் காட்டுக.
இதற்கு, !!a{surjective} சார்பு $f\colon \PosInt \to
A$ உள்ளது எனக் கொள்க.
!!a{surjective} சார்பு $g\colon \PosInt \to B$ ஐ வரையறுத்து,
அது !!{surjective} பண்புடையது என நிறுவுக.
$B = \emptyset$ எனில் என்ன நிகழ்கிறது?
\end{prob}

% SEGMENT OLP-0029-S16
\begin{prob}
$A_1$, $A_2$, \dots, $A_n$ ஆகிய அனைத்தும் !!{enumerable} எனில்,
$A_1 \cup \dots \cup A_n$ என்பதும் எண்ணிடத்தக்கது என்பதை,
$n$ மீதான தொகுத்தறிதல் மூலம் காட்டுக.
இரு கணங்கள் $A$, $B$ ஆகியவை !!{enumerable} எனில்,
$A \cup
B$ என்பதும் எண்ணிடத்தக்கது என்ற உண்மையை நீங்கள் முன்கொள்ளலாம்.
\end{prob}

% SEGMENT OLP-0029-S17
ஒரு கணத்தின் !!{element}s[acc] பட்டியலிடும்போது,
$1$ ஆவது !!{element} இலிருந்து எண்ணத் தொடங்குவது மேலும்
இயல்பாகத் தோன்றலாம். ஆயினும், கணிதவியலாளர்கள் பொருள்களை
எண்ணுவதற்கு இயல் எண்கள் $\Nat$ ஐப் பயன்படுத்த விரும்புகிறார்கள்.
அவர்கள் ஒரு பட்டியலின் $0$ ஆவது, $1$ ஆவது, $2$ ஆவது,
இவ்வாறான !!{element}s பற்றிப் பேசுகிறார்கள்.
அதற்கேற்ப, ஒரு பட்டியலாக்கத்தை $\Nat$ இலிருந்து $A$ க்குச் செல்லும்
!!a{surjective} சார்பாக வரையறுக்கலாம்.
இந்த இரண்டு வரையறைகளும் சமானமானவையே.

% SEGMENT OLP-0029-S18
\begin{prop}\ollabel{prop:enum-shift}
!!a{surjection} $g\colon \Nat \to A$ இருக்கும்போது,
அப்போதும் அப்போது மட்டுமே,
!!a{surjection} $f\colon \PosInt \to A$ உள்ளது.
\end{prop}

% SEGMENT OLP-0029-S19
\begin{proof}
!!a{surjection} $f\colon \PosInt \to A$ தரப்பட்டால்,
அனைத்து $n \in \Nat$ க்கும் $g(n) =
f(n+1)$ என வரையறுக்கலாம்.
$g\colon \Nat
\to A$ என்பது !!{surjective} பண்புடையது என்பதை எளிதில் காணலாம்.
மறுதிசையில், !!a{surjection} $g\colon
\Nat \to A$ தரப்பட்டால், $f(n) = g(n-1)$ என வரையறுக்கவும்.
\end{proof}

இது பின்வரும் முடிவைத் தருகிறது:

% SEGMENT OLP-0029-S20
\begin{cor}\ollabel{cor:enum-nat}
ஒரு கணம் $A$ வெற்றுக் கணமாக இருக்கும்போது அல்லது
!!a{surjective} சார்பு $f\colon \Nat \to A$ இருக்கும்போது,
அப்போதும் அப்போது மட்டுமே, அந்தக் கணம் !!{enumerable} ஆகும்.
\end{cor}

% SEGMENT OLP-0029-S21
ஒரு கணம் $A$ இன் !!{element}s கொண்ட பட்டியலை,
மீள இடம்பெறுதல்கள் இல்லாத பட்டியலாக மாற்ற முடியும் என்று
மேலே விவாதித்தோம். இது பட்டியலாக்கங்களுக்கும் பொருந்தும்;
ஆனால் இதைத் துல்லியமாகக் கூறி நிறுவுவது சற்றுக் கடினமானது.
எந்தச் சார்பு $f\colon \PosInt \to A$ உம்
அனைத்து $n \in
\PosInt$ க்கும் வரையறுக்கப்பட்டிருக்க வேண்டும்.
$A$ இல் முடிவுறு எண்ணிக்கையிலான !!{element}s மட்டுமே இருந்தால்,
$\PosInt$ இன் முடிவுறாத எண்ணிக்கையிலான !!{element}s மீது
வரையறுக்கப்பட்டு, $A$ இன் எல்லா !!{element}s[acc]யும் மதிப்புகளாக
எடுத்துக்கொண்டு, அதே மதிப்பை இருமுறை எடுக்காத சார்பு
இருக்க முடியாது என்பது தெளிவு.
அந்த நிலையில், அதாவது மீள இடம்பெறுதல்கள் இல்லாத பட்டியல்
முடிவுறுவதாக உள்ள நிலையில், $f$ க்கு வேறு சார்பகத்தைத்
தேர்ந்தெடுக்க வேண்டும்; அதில் முடிவுறு எண்ணிக்கையிலான
!!{element}s மட்டுமே இருக்க வேண்டும்.
மீள இடம்பெறுதல்கள் இல்லை என்பது $f$ ஆனது !!{injective}
பண்புடையதாக இருக்க வேண்டும் என்பதைக் குறிக்கிறது.
அது !!{surjective} பண்பையும் கொண்டிருப்பதால்,
ஏதேனும் ஒரு முடிவுறு கணம் $\{1, \dots, n\}$ அல்லது
$\PosInt$ என்பதற்கும் $A$ க்கும் இடையிலான
!!a{bijection} ஐத் தேடுகிறோம்.

% SEGMENT OLP-0029-S22
\begin{prop}\ollabel{prop:enum-bij}
$f\colon \PosInt \to A$ என்பது !!{surjective} பண்புடையது
(அதாவது $A$ இன் பட்டியலாக்கம்) எனில்,
!!a{bijection} $g\colon Z \to A$ உள்ளது.
இங்கு $Z$ என்பது $\PosInt$ அல்லது ஏதேனும்
$n \in \PosInt$ க்கு $\{1, \dots, n\}$ ஆகும்.
\end{prop}

% SEGMENT OLP-0029-S23
\begin{proof}
$g$ என்ற சார்பை மறுநிகழ்வு முறையில் வரையறுக்கிறோம்:
$g(1) = f(1)$ என அமைப்போம்.
$g(i)$ ஏற்கெனவே வரையறுக்கப்பட்டிருந்தால்,
$f(1)$, $f(2)$, \dots{} ஆகியவற்றில்,
$g(1)$, \dots, $g(i)$ ஆகியவற்றில் ஏற்கெனவே இடம்பெறாத
முதல் மதிப்பு ஏதாவது இருந்தால், அதை $g(i+1)$ என அமைப்போம்.
$A$ இல் சரியாக $n$ !!{element}s இருந்தால்,
$g(1)$, \dots, $g(n)$ ஆகிய அனைத்தும் வரையறுக்கப்பட்டுள்ளன;
எனவே $g\colon \{1, \dots, n\}
\to A$ என்ற சார்பை வரையறுத்துள்ளோம்.
$A$ இல் முடிவுறாத எண்ணிக்கையிலான !!{element}s இருந்தால்,
எந்த $i$ க்கும், $f(1)$, $f(2)$, \dots{} என்ற பட்டியலாக்கத்தில்,
$g(1)$, \dots, $g(i)$ ஆகியவற்றில் ஏற்கெனவே இடம்பெறாத
$A$ இன் !!a{element} இருக்க வேண்டும்.
இந்த நிலையில் $g\colon \PosInt \to A$ என்ற சார்பை வரையறுத்துள்ளோம்.

$g$ என்ற சார்பு !!{surjective} பண்புடையது.
ஏனெனில் $A$ இன் எந்த உறுப்பும் $f(1)$, $f(2)$, \dots{} ஆகியவற்றில்
இடம்பெறுகிறது ($f$ என்பது !!{surjective} பண்புடையதால்);
எனவே இறுதியில் ஏதாவது ஒரு $i$ க்கு $g(i)$ இன் மதிப்பாக இருக்கும்.
அது !!{injective} பண்பையும் கொண்டுள்ளது.
ஏனெனில் $g(j) = g(i)$ என அமையும் $j < i$ இருந்தால்,
$g(i)$ என்பது $g(1)$, \dots, $g(i-1)$ ஆகியவற்றில்
ஏற்கெனவே இடம்பெற்றிருக்கும்.
இது நாம் $g$ ஐ வரையறுத்த விதத்திற்கு முரணாகும்.
\end{proof}

% SEGMENT OLP-0029-S24
\begin{cor}\ollabel{cor:enum-nat-bij}
ஒரு கணம் $A$ வெற்றுக் கணமாக இருக்கும்போது அல்லது
!!a{bijection} $f\colon N \to A$ இருக்கும்போது,
அப்போதும் அப்போது மட்டுமே, அது !!{enumerable} ஆகும்.
இங்கு $N = \Nat$ அல்லது ஏதேனும் $n \in \Nat$ க்கு
$N = \{0, \dots, n\}$ ஆகும்.
\end{cor}

% SEGMENT OLP-0029-S25
\begin{proof}
$A$ வெற்றுக் கணமாக இருக்கும்போது அல்லது
!!a{surjective} சார்பு $f\colon \PosInt \to A$ இருக்கும்போது,
அப்போதும் அப்போது மட்டுமே, $A$ என்பது !!{enumerable} ஆகும்.
\olref{prop:enum-bij} இன்படி, பிந்தைய நிபந்தனை நிறைவேறுவதற்குத்
தேவையானதும் போதுமானதுமான நிபந்தனை,
!!a{bijective} சார்பு $f\colon Z \to A$ இருப்பதாகும்.
இங்கு $Z =
\PosInt$ அல்லது ஏதேனும் $n \in \PosInt$ க்கு
$Z = \{1, \dots, n\}$ ஆகும்.
\olref{prop:enum-shift} இன் நிறுவலில் பயன்படுத்திய அதே வாதத்தின்படி,
இது நிறைவேறுவதற்குத் தேவையானதும் போதுமானதுமான நிபந்தனை,
!!a{bijection} $g\colon N \to A$ இருப்பதாகும்;
இங்கு $N = \Nat$ அல்லது $N = \{0, \dots, n-1\}$ ஆகும்.
\end{proof}

% SEGMENT OLP-0029-S26
\begin{prob}
\olref[sfr][siz][enm]{defn:enumerable} இன்படி,
$A = \emptyset$ என இருக்கும்போது அல்லது
!!a{surjective} சார்பு $f\colon
\PosInt \to A$ இருக்கும்போது, அப்போதும் அப்போது மட்டுமே,
ஒரு கணம் $A$ எண்ணிடத்தக்கதாகும்.
“!!{enumerable} கணம்” என்பதைத் துல்லியமாகப் பின்வருமாறும்
வரையறுக்கலாம்: !!a{injective} சார்பு
$g\colon A \to \PosInt$ இருக்கும்போது,
அப்போதும் அப்போது மட்டுமே, ஒரு கணம் எண்ணிடத்தக்கதாகும்.
இந்த வரையறைகள் சமானமானவை எனக் காட்டுக.
அதாவது, $A = \emptyset$ என இருக்கும்போது அல்லது
!!a{surjective} சார்பு $f\colon \PosInt \to A$ இருக்கும்போது,
அப்போதும் அப்போது மட்டுமே,
!!a{injective} சார்பு $g\colon A \to \PosInt$ உள்ளது எனக் காட்டுக.
\end{prob}

\end{document}
